\title{DeepSeekMath: Pushing the Limits of Mathematical Reasoning in Open Language Models}

\begin{abstract}
Mathematical reasoning remains a formidable obstacle for language models due to its inherently intricate and highly organized structure. In this study, we present DeepSeekMath~7B, which is built upon DeepSeek-Coder-Base-v1.5 7B through continued pre-training with a dataset of 120B math-focused tokens collected from Common Crawl, supplemented by both natural language and programming code. DeepSeekMath~7B attains a noteworthy 51.7\% accuracy on the MATH benchmark at the competition level, without dependence on external toolkits or ensemble voting strategies, nearing the performance exhibited by Gemini-Ultra and GPT-4. When leveraging self-consistency over 64 generated samples, DeepSeekMath~7B further raises its MATH score to 60.9\%. The enhanced mathematical reasoning capabilities of DeepSeekMath~7B can be traced to two principal innovations: firstly, the exploitation of large-scale web data via a carefully designed data selection framework; secondly, the introduction of Group Relative Policy Optimization (GRPO), a Proximal Policy Optimization (PPO) variant, which both strengthens mathematical reasoning and improves PPO memory efficiency.
\end{abstract}

\section{Introduction}

Large language models (LLMs) have transformed methodologies for mathematical reasoning within the field of artificial intelligence, leading to notable progress on benchmarks for quantitative reasoning \citep{MATH} as well as geometry reasoning \citep{trinh2024solving}. In addition, these models have demonstrated value in supporting humans with the resolution of intricate mathematical tasks \citep{tao}. Nonetheless, state-of-the-art models like GPT-4 \citep{gpt4} and Gemini-Ultra \citep{gemini} remain inaccessible to the public, and existing open-source alternatives lag significantly in terms of performance.

This paper presents DeepSeekMath, a specialized language model engineered for mathematics that demonstrates superior performance compared to existing open-source models and narrows the gap to GPT-4 across standard academic assessments. Central to this advance is the DeepSeekMath Corpus, a massive and meticulously curated pre-training dataset containing 120B math-related tokens. The corpus was assembled from the Common Crawl (CC) using a classifier built on fastText \citep{joulin2016fasttext}. Initially, training data for the classifier was sourced from OpenWebMath \citep{openwebmath} for positive samples, with a varied assortment of non-mathematical web content constituting negative samples. The classifier was then used to identify additional mathematical content within the CC, which underwent further selection and validation by human annotators. These validated samples were leveraged to iteratively improve the classifier’s accuracy. Empirical assessments confirm the corpus’s high caliber, as evidenced by DeepSeekMath-Base 7B’s attainment of 64.2\% on GSM8K \citep{gsm8k} and 36.2\% on the MATH competition dataset \citep{MATH}, surpassing the results achieved by Minerva 540B \citep{minerva}. Moreover, thanks to its multilingual construction, DeepSeekMath exhibits enhanced outcomes on Chinese mathematics evaluation tasks \citep{wei2023cmath,agieval}. The methodology and insights gained in the compilation of mathematical training data presented here may serve as a foundational resource for further progress within the research community, with ample opportunities for future enhancements.

DeepSeekMath-Base is built upon DeepSeek-Coder-Base-v1.5 7B \citep{deepseek-coder}, based on our findings that initializing from a model pretrained on code yields superior outcomes to starting with a generic LLM. In addition, mathematical pre-training leads to enhancements not only in mathematical tasks but also in broader reasoning, as evidenced by improved performance on both MMLU \citep{mmlu} and BBH \citep{bbh} benchmarks.

Subsequent to pre-training, we perform mathematical instruction fine-tuning on DeepSeekMath-Base, utilizing data from chain-of-thought \citep{cot}, program-of-thought \citep{pot,pal}, and tool-integrated reasoning \citep{tora}. The final model, DeepSeekMath-Instruct 7B, surpasses all other 7B-sized peers and rivals open-source instruction-tuned models with 70B parameters.

In addition, we propose Group Relative Policy Optimization (GRPO), a novel reinforcement learning (RL) algorithm derived from Proximal Policy Optimization (PPO) \citep{schulman2017proximal}. Unlike PPO, GRPO eliminates the need for a critic network, instead using group-based score estimations to establish the baseline, thereby considerably lowering computational demands during training. Utilizing only a selection of English instruction-tuning data, GRPO demonstrates marked gains over the robust DeepSeekMath-Instruct model, achieving higher scores in both in-domain tasks (GSM8K: 82.9\% $\rightarrow$ 88.2\%, MATH: 46.8\% $\rightarrow$ 51.7\%) and out-of-domain assessments (e.g., CMATH: 84.6\% $\rightarrow$ 88.8\%) in the reinforcement learning stage. Furthermore, we articulate a unified framework that interprets a range of methods—Rejection Sampling Fine-Tuning (RFT) \citep{yuan2023scaling}, Direct Preference Optimization (DPO) \citep{dpo}, PPO, and GRPO—under a shared conceptual lens, identifying them all as either direct or simplified approaches within RL. To further elucidate this framework, we carry out comprehensive studies, comparing, for example, online versus offline learning, process versus outcome supervision, and single-turn versus iterative RL, thus dissecting the critical facets of the paradigm. Finally, we analyze the mechanisms through which our RL technique enhances the effectiveness of instruction-tuned models, and outline promising future directions for advancing RL methodologies under this unifying framework.

\subsection{Contributions}
Our work offers scalable math pre-training as well as a thorough examination and evaluation of reinforcement learning methods.

\textbf{Math Pre-Training at Scale}

\begin{itemize}[topsep=0pt]
    \item Our study demonstrates that the openly available Common Crawl dataset holds significant utility for mathematical applications. Through the implementation of a carefully crafted data filtering and selection pipeline, we assemble the DeepSeekMath Corpus—an extensive, high-quality dataset comprising 120B tokens derived from web sources specifically curated for mathematical relevance. This dataset is nearly 7 times larger than the collection of math web pages utilized in Minerva \citep{minerva}, and 9 times larger than the recently introduced OpenWebMath \citep{openwebmath}.

    \item The DeepSeekMath-Base 7B model, upon pre-training, delivers performance on par with Minerva 540B \citep{minerva}, suggesting that model size alone does not solely determine effectiveness in mathematical reasoning tasks. In fact, even models with fewer parameters, when trained on data of exceptional quality, are capable of exhibiting robust performance.
\end{itemize}

\item We present the results from our experiments on mathematical training. Pre-training models with code enhances their proficiency in tackling mathematical problems, regardless of whether auxiliary tools are involved. This provides some insight into the debated issue: \textit{can code training bolster reasoning capabilities?} Our results indicate that it does, specifically in the context of mathematical reasoning.

    \item Despite the widespread practice of incorporating arXiv papers into training—particularly in mathematics-focused studies—our analysis finds no significant gains across any of the mathematical benchmarks considered in this work.
\end{itemize}

\textbf{Exploration and Analysis of Reinforcement Learning}

\begin{itemize}[topsep=0pt]
    \item We present Group Relative Policy Optimization (GRPO), a reinforcement learning technique that is both resource-efficient and highly effective. Unlike traditional approaches, GRPO eliminates the need for a critic model by deriving the baseline from group-level scores, thereby substantially lowering the computational demands relative to Proximal Policy Optimization (PPO).
    \item Our results show that GRPO delivers substantial improvements to the performance of our instruction-tuned system DeepSeekMath-Instruct using only the instruction-tuning dataset. In addition, reinforcement learning with GRPO leads to gains in out-of-domain generalization.
    \item We propose a comprehensive framework that offers a unified perspective on several methodologies, including RFT, DPO, PPO, and GRPO. We further perform in-depth empirical studies encompassing comparisons such as online versus offline learning, outcome-based versus process-based supervision, and single-turn versus iterative reinforcement learning, to analyze the fundamental aspects of this unified paradigm.
    \item Leveraging this unified perspective, we examine why reinforcement learning is effective, and outline promising avenues for further improving the application of reinforcement learning techniques to LLMs.
\end{itemize}

\subsection{Summary of Evaluations and Metrics}

\begin{itemize}[topsep=0pt]
    \item \textbf{English and Chinese Mathematical Reasoning}: We perform thorough evaluations of our models utilizing both English and Chinese benchmark datasets, spanning mathematical tasks from elementary through university levels. The English evaluations consist of GSM8K \citep{gsm8k}, MATH \citep{MATH}, SAT \citep{llemma}, OCW Courses \citep{minerva}, and MMLU-STEM \citep{mmlu}. For Chinese, we employ MGSM-zh \citep{mgsm}, CMATH \citep{wei2023cmath}, Gaokao-MathCloze \citep{agieval}, and Gaokao-MathQA \citep{agieval}. Our analysis considers both the capability of models to independently produce complete textual solutions and their proficiency in problem-solving via Python.

DeepSeekMath-Base demonstrates comparable results to the proprietary Minerva 540B model \citep{minerva} on English evaluation tasks and outperforms all available open-source base models, such as Mistral 7B \citep{mistral} and Llemma-34B \citep{llemma}, independent of whether those models were specifically pre-trained on mathematical data—often by a substantial degree. Importantly, DeepSeekMath-Base achieves even higher scores on Chinese benchmarks. This improvement is likely attributed to our decision to broaden the pre-training corpus beyond English-only mathematical data, unlike previous studies \citep{minerva,llemma}, and to incorporate extensive high-quality math datasets in other languages. Further advancements through mathematical instruction tuning and reinforcement learning yield DeepSeekMath-Instruct and DeepSeekMath-RL, which both achieve outstanding accuracy, surpassing 50\% on the challenging MATH dataset—marking a first within the open-source community.

\item \textbf{Formal Mathematics}: DeepSeekMath-Base is assessed on the informal-to-formal theorem proving benchmark from \citep{dsp_proof}, employing the miniF2F dataset \citep{minif2f} and utilizing Isabelle \citep{isabelle} as the proof assistant. The model exhibits robust few-shot autoformalization capabilities.
\item \textbf{Natural Language Understanding, Reasoning, and Code}: To thoroughly gauge the general proficiency of the model in understanding, reasoning, and programming, we conduct evaluations of DeepSeekMath-Base using the Massive Multitask Language Understanding (MMLU) benchmark \citep{mmlu}, which includes 57 multiple-choice problems spanning a wide range of disciplines. Additionally, we use BIG-Bench Hard (BBH) \citep{bbh}, composed of 23 tasks demanding predominantly multi-step reasoning, and also assess performance on HumanEval \citep{codex} and MBPP \citep{mbpp}, both of which are standard benchmarks for code-oriented models. Pre-training on mathematics data leads to measurable improvements in both language comprehension and reasoning abilities.

\section{Math Pre-Training}

\subsection{Data Collection and Decontamination} This section details the methodology used to assemble the DeepSeekMath Corpus using data sourced from Common Crawl. Figure \ref{fig:data_collect} illustrates the stepwise pipeline implemented for systematically extracting a large-scale mathematics dataset, initiated by leveraging a seed corpus—typically a concise yet high-quality math-oriented dataset. Notably, this methodology is readily transferable to additional domains, including programming.

Initially, we utilize OpenWebMath \citep{openwebmath}, a repository of high-quality mathematical texts sourced from the web, as our primary seed dataset. Leveraging this seed corpus, we train a fastText model \citep{joulin2016fasttext} to retrieve additional web pages exhibiting similar mathematical content to OpenWebMath. Concretely, we designate 500,000 randomly chosen entries from the seed corpus as positive samples, while 500,000 web pages drawn from Common Crawl serve as negative samples. For model training, we adopt an open-source fastText implementation\footnote{\url{https://fasttext.cc}}, configuring it with a 256-dimensional embedding space, a learning rate of 0.1, a maximum word n-gram length of 3, a minimum threshold of 3 word occurrences, and three training epochs. To efficiently handle the vast Common Crawl dataset, we implement deduplication at the URL level as well as near-deduplication, which reduces the dataset to 40B unique HTML web pages. The trained fastText model is then used to identify and select mathematical content from this deduplicated set. To eliminate low-quality material, we sort the retrieved pages based on their predicted fastText scores, retaining only the highest-scoring documents. The adequacy of the selected data volume is examined by conducting pre-training experiments on subsets containing the top 40B, 80B, 120B, and 160B tokens. For the initial round, we opt to preserve those pages comprising the top 40B tokens.

Following the initial round of data acquisition, a considerable number of mathematical web pages remain uncaptured, primarily because the fastText model's set of positive examples does not exhibit adequate variety. To address this limitation, we seek out further sources of mathematical web content to expand the seed corpus and, in turn, enhance the performance of the fastText model. Our process begins by partitioning the entirety of Common Crawl into non-overlapping domains, where each domain comprises web pages that share the same base URL. For each domain, we determine the proportion of web pages retrieved during the initial round. Domains in which more than 10\% of the pages were acquired are designated as math-related domains (such as \url{mathoverflow.net}).

We then proceed to manually tag URLs within these selected domains that pertain to mathematical subject matter (for instance, \url{mathoverflow.net/questions}). Any web pages associated with these URLs that were not previously collected are incorporated into the seed corpus. By broadening the corpus of positive examples in this way, the next iteration of the fastText model is trained to better recognize and extract mathematical material. Repeating this procedure through four rounds of data gathering results in a total of 35.5M mathematical web pages and an aggregate of 120B tokens. During the fourth cycle, it becomes apparent that almost 98\% of the documents had already been gathered by the third iteration, prompting us to halt the data collection process at that stage.

In order to prevent contamination of benchmarks, we adopt the approach described in \cite{deepseek-coder} to exclude web content featuring questions or answers from English mathematical benchmarks like GSM8K \citep{gsm8k} and MATH \citep{MATH}, as well as Chinese benchmarks such as CMATH \citep{wei2023cmath} and AGIEval \citep{agieval}. Specifically, our filtering process eliminates any segment of text containing a 10-gram sequence that exactly corresponds to any substring found in these evaluation benchmarks from our mathematical training set. For benchmark texts whose length falls below 10 grams but comprises at least 3 grams, we rely on exact match filtering to remove any potentially contaminated web data.

\subsection{Validating the Quality of the DeepSeekMath~Corpus}

To evaluate the DeepSeekMath~Corpus in relation to other recently introduced mathematics-focused training datasets, we conduct pre-training experiments using the following corpora:

\begin{itemize}[topsep=0pt]
    \item \textbf{MathPile} \citep{mathpile}: an 8.9B token multi-source dataset, compiled from materials such as textbooks, Wikipedia, ProofWiki, CommonCrawl, StackExchange, and predominantly arXiv (over 85\% of the content);
    \item \textbf{OpenWebMath} \citep{openwebmath}: a 13.6B token collection derived by filtering CommonCrawl for web data related to mathematics;
    \item \textbf{Proof-Pile-2} \citep{llemma}: a comprehensive mathematical dataset that includes OpenWebMath, AlgebraicStack (comprising 10.3B tokens of mathematical code), and 28.0B tokens of arXiv documents. In accordance with \cite{llemma}, we apply an arXiv:Web:Code distribution ratio of 2:4:1 for experiments utilizing Proof-Pile-2.
\end{itemize}

\subsubsection{Training Setting}
\label{sec:quality-policy}
We conduct math-specific training on a general-purpose language model containing 1.3B parameters, which adopts the architecture utilized by DeepSeek LLMs \citep{deepseek-llm}, and refer to this as DeepSeek-LLM 1.3B. Each mathematical dataset is used to independently train a model over 150B tokens. For all experimental runs, we leverage the HAI-LLM \citep{haillm} platform, recognized for its efficiency and lightweight design. Adhering to the established training regimen of DeepSeek LLMs, we employ the AdamW optimizer \citep{adamW} with hyperparameters $\beta_1=0.9$, $\beta_2=0.95$, and a weight decay of $0.1$. The learning rate is managed with a multi-phase schedule: it ascends to its maximum after 2{,}000 warm-up iterations, falls to 31.6\% of the peak following 80\% completion of training, and is reduced to 10.0\% at 90\% completion. We cap the peak learning rate at $5.3\text{e-}4$, use a batch size of 4M tokens, and fix the context window at 4K tokens.

\subsubsection{Evaluation Results}

\textbf{The DeepSeekMath~Corpus is of high quality, covers multilingual mathematical content, and is the largest in size.}

\begin{itemize}[topsep=0pt]
    \item \textbf{High-quality}: We assess downstream capabilities across eight mathematical benchmarks using few-shot chain-of-thought prompting \cite{cot}. According to the results in Table \ref{tab:corpora_comparison}, models trained with the DeepSeekMath~Corpus consistently achieve superior outcomes. Furthermore, Figure \ref{fig:corpora_comparisons} illustrates that at the 50B token mark (equivalent to one complete epoch of Proof-Pile-2), the DeepSeekMath-trained model outperforms the one trained on Proof-Pile-2, underscoring the generally higher quality of the DeepSeekMath Corpus.
    \item \textbf{Multilingual}: The DeepSeekMath~Corpus includes a diverse range of languages, with English and Chinese forming the largest portions of the dataset. As displayed in Table \ref{tab:corpora_comparison}, leveraging DeepSeekMath~Corpus in training significantly boosts mathematical reasoning performance for both English and Chinese. In comparison, previously available mathematical corpora are heavily dominated by English and tend to yield limited or even detrimental effects on Chinese mathematical reasoning tasks.
    \item \textbf{Large-scale}: DeepSeekMath~Corpus is substantially larger than other mathematical datasets presently available. Figure \ref{fig:corpora_comparisons} demonstrates that, when trained on DeepSeekMath~Corpus, DeepSeek-LLM 1.3B exhibits a sharper ascent in performance and more sustained improvements. Conversely, the alternative corpora are far smaller and undergo repeated passes during training, resulting in model progress quickly stagnating.
\end{itemize}

\subsection{Training and Evaluating DeepSeekMath-Base 7B}

This section presents DeepSeekMath-Base 7B, a foundational model exhibiting advanced mathematical reasoning skills. Built upon DeepSeek-Coder-Base-v1.5 7B \citep{deepseek-coder}, the model underwent training on a corpus totaling 500B tokens. The data utilized comprised 56\% DeepSeekMath Corpus, 4\% AlgebraicStack, 10\% arXiv, 20\% code from Github, and the remaining 10\% consisted of natural language examples sourced from Common Crawl in both English and Chinese. While the majority of the training configurations follow those outlined in Section \ref{sec:quality-policy}, we increase the peak learning rate to 4.2e-4 and employ a batch size of 10M tokens.

In this study, we perform an extensive evaluation of the mathematical proficiency of DeepSeekMath-Base 7B. Our analysis centers on the model's capacity to independently generate complete mathematical solutions, its effectiveness in solving mathematical problems with the assistance of tools, and its aptitude for formal theorem proving. Additionally, we extend our investigation beyond mathematics to present a broader characterization of the base model, detailing its abilities in natural language understanding, reasoning, and programming.

\paragraph{Mathematical Problem Solving with Step-by-Step Reasoning}

We assess the capability of DeepSeekMath-Base to tackle mathematical problems utilizing few-shot chain-of-thought prompting \citep{cot}, evaluating its performance on eight distinct benchmarks in both English and Chinese. These benchmarks span quantitative reasoning tasks—such as GSM8K \citep{gsm8k}, MATH \citep{MATH}, and CMATH \citep{wei2023cmath}—as well as multiple-choice assessments like MMLU-STEM \citep{mmlu} and Gaokao-MathQA \citep{agieval}. Together, they encompass a broad spectrum of mathematical domains, ranging from elementary mathematics up to college-level topics.

As indicated in Table \ref{tab:base_math_cot}, DeepSeekMath-Base 7B demonstrates superior results across all eight benchmarks when compared with other open-source base models—including both the broadly adopted general-purpose Mistral 7B \citep{mistral} and the recently introduced Llemma 34B \citep{llemma}, the latter of which is specifically trained on Proof-Pile-2 \citep{llemma} for mathematical tasks. Particularly remarkable is DeepSeekMath-Base’s performance on the MATH dataset, where it exceeds all available open-source base models by more than 10\% in absolute terms and even outperforms Minerva 540B \citep{minerva}, a substantially larger closed-source model (77 times greater in parameter count), which itself is based on PaLM \citep{palm} and has undergone extensive additional training with mathematical content.

\paragraph{Mathematical Problem Solving with Tool Use} We assess the capability of models to perform program-assisted mathematical reasoning on the GSM8K and MATH benchmarks, leveraging few-shot program-of-thought prompting \citep{pot,pal}. The models are instructed to construct Python scripts to address each problem, making use of libraries like \textit{math} and \textit{sympy} to handle complex calculations. The solutions are judged based on the output produced by running the generated program. As indicated in Table \ref{tab:base_math_pot_proof}, DeepSeekMath-Base 7B achieves superior results compared to the previous leading model, Llemma 34B.

\paragraph{Formal Mathematics}

The automation of formal proof has gained substantial traction recently, as it contributes significantly to verifying the precision and dependability of mathematical proofs while also improving efficiency. In this context, we assess DeepSeekMath-Base 7B on the informal-to-formal proving task introduced in \citep{dsp_proof}, which involves producing a formal proof from an informal statement, its formal equivalent, and an informal proof. Our experiments utilize miniF2F \citep{minif2f}, a benchmark designed for formalizing Olympiad-level mathematics. For each problem, we prompt the model with a few examples to generate a formal proof in Isabelle. Mirroring the methodology in \cite{dsp_proof}, we have the model generate proof outlines, after which we employ the standard automated prover Sledgehammer \citep{sledgehammer} to complete any missing reasoning steps. As indicated in Table \ref{tab:base_math_pot_proof}, DeepSeekMath-Base 7B achieves robust results in the automatic formalization of proofs.

\paragraph{Natural Language Understanding, Reasoning, and Code}

We assess natural language understanding using MMLU \citep{mmlu}, evaluate reasoning with BBH \citep{bbh}, and examine programming abilities via HumanEval \citep{codex} and MBPP \citep{mbpp}. Table \ref{tab:understanding_reasoning_code} highlights that DeepSeekMath-Base 7B achieves notable improvements over its predecessor, DeepSeek-Coder-Base-v1.5 \citep{deepseek-coder}, on both MMLU and BBH tasks. These results underscore the beneficial influence of mathematical training on language understanding and reasoning. Furthermore, by incorporating code tokens in ongoing training, DeepSeekMath-Base 7B retains the coding performance of DeepSeek-Coder-Base-v1.5 across both programming benchmarks. Taken together, DeepSeekMath-Base 7B consistently surpasses the general-purpose Mistral 7B model \citep{mistral} on all three benchmarks related to reasoning and coding.

\section{Supervised Fine-Tuning}

\subsection{SFT Data Curation}

We develop a comprehensive mathematical instruction-tuning dataset that encompasses both English and Chinese questions, spanning a range of mathematical domains and difficulty levels. Each problem is coupled with detailed solutions formatted as chain-of-thought (CoT) \citep{cot}, program-of-thought (PoT) \citep{pot,pal}, or tool-integrated reasoning \citep{tora}. In total, our training dataset comprises 776K examples.

\begin{itemize}[topsep=0pt]
    \item \textbf{English mathematical datasets}: Our dataset includes tool-integrated solution annotations for GSM8K and MATH problems. Additionally, we incorporate a selected subset of MathInstruct \citep{MathInstruct} as well as the training split from Lila-OOD \citep{lila}, where problems are addressed using CoT or PoT methodologies. The English dataset features a broad spectrum of mathematical topics, including algebra, probability, number theory, calculus, and geometry.
    \item \textbf{Chinese mathematical datasets}: We gather K-12 mathematics problems in Chinese, distributed over 76 sub-categories such as linear equations, and provide corresponding solutions using both CoT and tool-integrated reasoning approaches.
\end{itemize}

\subsection{Training and Evaluating DeepSeekMath-Instruct 7B}

In this section, we present DeepSeekMath-Instruct 7B, which is developed by applying mathematical instruction tuning to DeepSeekMath-Base. During training, examples are concatenated in random order until the input sequence attains a maximum context length of 4K tokens. The model is trained over 500 iterations, utilizing a batch size of 256 and maintaining a fixed learning rate of 5e-5.

To assess mathematical capabilities, we evaluate the models under both tool-free and tool-assisted conditions, utilizing four quantitative reasoning benchmarks available in both English and Chinese. Our evaluation includes comparisons with the top-performing models available at the time:

\begin{itemize}[topsep=0pt]
    \item \textbf{Closed-source models} consist of: (1) the GPT series, with GPT-4 \citep{gpt4} and the GPT-4 Code Interpreter~\footnote{\url{https://openai.com/blog/chatgpt-plugins\#\#code-interpreter}} representing the most advanced variants, (2) Gemini Ultra and Pro \citep{gemini}, (3) Inflection-2 \citep{inflection-2}, (4) Grok-1~\footnote{\url{https://x.ai/model-card}}, along with several models introduced by Chinese enterprises, such as (5) Baichuan-3~\footnote{\url{https://www.baichuan-ai.com}}, and (6) the most recent addition to the GLM line, GLM-4~\footnote{\url{https://open.bigmodel.cn/dev/api\#glm-4}} from the GLM family \citep{glm}.

These models are designed for broad applicability, with most having completed various stages of alignment optimization.
    \item \textbf{Open-source models} encompass: generic language models such as (1) DeepSeek-LLM-Chat 67B \citep{deepseek-llm}, (2) Qwen 72B \citep{qwen}, (3) SeaLLM-v2 7B \citep{seallm}, and (4) ChatGLM3 6B \citep{chatglm3}. Additionally, there are models specialized in mathematical capabilities, including (5) InternLM2-Math 20B~\footnote{\url{https://github.com/InternLM/InternLM-Math}}, derived from InternLM2, which received targeted mathematical pre-training followed by instruction tuning, (6) Math-Shepherd-Mistral 7B, which enhances Mistral 7B \citep{mistral} through PPO training \citep{schulman2017proximal} with supervision from a process-oriented reward model, (7) the WizardMath collection \citep{wizardmath}, which advances math reasoning in both Mistral 7B and Llama-2 70B \citep{llama2} via evolve-instruct—an instruction tuning methodology utilizing AI-generated prompts—and PPO, with training datasets predominantly from GSM8K and MATH, (8) MetaMath 70B \citep{metamath}, an adaptation of Llama-2 70B fine-tuned using a data-augmented blend of GSM8K and MATH, (9) ToRA 34B \cite{tora}, a version of CodeLlama 34B that has been specifically trained for tool-supported mathematical reasoning, and (10) MAmmoTH 70B \citep{MathInstruct}, which is an instruction-tuned variant of Llama-2 70B employing MathInstruct as its supervision.

According to Table \ref{tab:sft_rl_math}, when evaluated without the aid of external tools, DeepSeekMath-Instruct 7B excels in executing detailed, stepwise reasoning. Specifically, on the competition-level MATH benchmark, this model outperforms all publicly available models and most proprietary alternatives (such as Inflection-2 and Gemini Pro) by a margin of at least 9\% in absolute terms. This advantage persists even when compared to models with significantly larger parameter counts (such as Qwen 72B) or those that have received additional reinforcement learning tailored to mathematics (e.g., WizardMath-v1.1 7B). While DeepSeekMath-Instruct competes closely with proprietary Chinese models like GLM-4 and Baichuan-3 on MATH, it still lags behind leading systems such as GPT-4 and Gemini Ultra.

When evaluated under conditions that permit the use of both natural language reasoning and programmatic tool deployment for solving problems, DeepSeekMath-Instruct 7B attains nearly 60\% accuracy on the MATH benchmark, outperforming all previously released open-source alternatives. Furthermore, across additional benchmarks, our model demonstrates performance on par with DeepSeek-LLM-Chat 67B, the previous state-of-the-art model which is an order of magnitude larger.

\section{Reinforcement Learning}

\subsection{Group Relative Policy Optimization} Previous research has shown that reinforcement learning (RL) can further enhance the mathematical reasoning skills of large language models (LLMs) after they have undergone Supervised Fine-Tuning (SFT) \citep{wang2023math,wizardmath}. Here, we present our novel and resource-efficient RL approach, Group Relative Policy Optimization (GRPO).

\subsubsection{From PPO to GRPO} Proximal Policy Optimization (PPO) \citep{schulman2017proximal} is a commonly adopted actor-critic reinforcement learning approach, particularly prevalent in the RL-based fine-tuning phase of LLM development \citep{ouyang2022training}. The central idea behind PPO is to update LLMs by maximizing the following clipped surrogate objective function:

\begin{equation}
\footnotesize
    \mathcal{J}_{PPO}(\theta) = \mathbb{E}{[q \sim P(Q), o \sim \pi_{\theta_{old}}(O|q)]} \frac{1}{|o|} \sum_{t=1}^{|o|} \min \left[ \frac{\pi_\theta(o_{t} | q, o_{<t})}{\pi_{\theta_{old}}(o_{t} | q, o_{<t})} A_{t}, \text{clip} \left( \frac{\pi_\theta(o_{t} | q, o_{<t})}{\pi_{\theta_{old}}(o_{t} | q, o_{<t})}, 1 - \varepsilon, 1 + \varepsilon \right)  A_{t} \right] ,
\end{equation}

In this context, $\pi_{\theta}$ and $\pi_{\theta_{old}}$ denote the present and previous policy models, respectively, while $q$ and $o$ refer to questions and outputs that are drawn from the question dataset and the prior policy $\pi_{\theta_{old}}$. The hyper-parameter $\varepsilon$, associated with clipping, is employed in PPO to promote stable training dynamics. The term $A_t$ represents the advantage, calculated using Generalized Advantage Estimation (GAE) \citep{gae}, which leverages rewards $\{r_{\ge t}\}$ together with a trainable value function $V_{\psi}$. Consequently, PPO requires the simultaneous training of a value function in addition to the policy model. To counteract excessive reward model optimization, it is customary to integrate a per-token KL divergence penalty against a reference model into the reward for each token, as suggested by \citet{ouyang2022training}, formally given by

\begin{equation}
    r_{t} = r_\varphi(q, o_{\le t}) - \beta \log\frac{\pi_{\theta}(o_{t}|q, o_{<t})}{\pi_{ref}(o_{t}|q, o_{<t})},
\label{eq:PPO-reward}
\end{equation}

Here, $r_\varphi$ denotes the reward model, while $\pi_{ref}$ stands for the reference model—most often the initial SFT model—and $\beta$ represents the scaling factor for the KL penalty term.

Since the value function utilized in PPO generally constitutes a separate model of a size similar to the policy network, it imposes considerable memory and computation demands. Moreover, within the reinforcement learning paradigm, the value function serves as a baseline for computing the advantage, thus helping to reduce variance. However, in the setting of LLMs, the reward model typically provides a reward only for the final token, making it more challenging to train a value function that remains precise across all tokens. To overcome these issues, as illustrated in Figure \ref{fig:grpo}, we introduce Group Relative Policy Optimization (GRPO). Unlike PPO, GRPO eliminates the requirement for a distinct value function approximation and instead relies on the average reward obtained from several outputs, each generated for the same input prompt, to serve as a baseline. Concretely, for any given question $q$, GRPO generates a collection of responses $\{o_1, o_2, \cdots, o_G\}$ by sampling from the prior policy $\pi_{\theta_{old}}$. The policy model is then refined by optimizing the following objective:

\begin{equation}
\footnotesize
\begin{split}
    \mathcal{J}_{GRPO}(\theta) &= \mathbb{E}{[q \sim P(Q), \{o_i\}_{i=1}^G \sim \pi_{\theta_{old}}(O|q)]}  \\
    & \frac{1}{G}\sum_{i=1}^G\frac{1}{|o_i|} \sum_{t=1}^{|o_i|} \left\{ \min \left[ \frac{\pi_\theta(o_{i,t} | q, o_{i,<t})}{\pi_{\theta_{old}}(o_{i,t} | q, o_{i,<t})} \hat{A}_{i,t}, \text{clip} \left( \frac{\pi_\theta(o_{i,t} | q, o_{i,<t})}{\pi_{\theta_{old}}(o_{i,t} | q, o_{i,<t})}, 1 - \varepsilon, 1 + \varepsilon \right)  \hat{A}_{i,t} \right] - \beta \mathbb{D}_{KL}\left[\pi_{\theta} || \pi_{ref}\right]\right\} ,
\end{split}
\label{eq:GRPO-obj}
\end{equation}

where $\varepsilon$ and $\beta$ serve as hyper-parameters, and $\hat{A}_{i,t}$ denotes the advantage computed solely from the relative rewards of outputs within the same group, with further explanation provided in subsequent subsections. The method by which GRPO computes advantages—using only group-relative rewards—is particularly well-suited to the inherently comparative approach of reward models, as such models are generally trained on collections of pairwise output comparisons for each prompt. It is also important to observe that, rather than incorporating a KL penalty into the reward itself, GRPO instead introduces regularization by directly adding the KL divergence between the policy under training and the reference policy to the objective, thereby sidestepping additional complexity in calculating $\hat{A}_{i,t}$. In contrast to the KL penalty adopted in (\ref{eq:PPO-reward}), our KL divergence is computed using the following unbiased estimator \citep{kl_approx}:

\begin{equation}
\small
    \mathbb{D}_{KL}\left[\pi_{\theta} || \pi_{ref}\right] = \frac{\pi_{ref}(o_{i,t}|q,o_{i,<t})}{\pi_{\theta}(o_{i,t}|q,o_{i,<t})}- \log\frac{\pi_{ref}(o_{i,t}|q,o_{i,<t})}{\pi_{\theta}(o_{i,t}|q,o_{i,<t})} - 1,
\end{equation}

which is always non-negative by construction.

\begin{algorithm}[t]
  \small
  \caption{Iterative Group Relative Policy Optimization}
  \textbf{Input} initial policy model $\pi_{\theta_{\text{init}}}$; reward models $r_\varphi$; task prompts $\mathcal{D}$; 
  hyperparameters $\varepsilon$, $\beta$, $\mu$
  \begin{algorithmic}[1]
    \State Set policy model $\pi_\theta \gets \pi_{\theta_{\text{init}}}$
    \For{iteration = 1, \ldots, I}
       \State Assign current policy as reference model: $\pi_{ref} \gets \pi_{\theta}$
      \For{step = 1, \ldots, M}
      \State Draw a minibatch $\mathcal{D}_b$ from $\mathcal{D}$
      \State Save current policy model $\pi_{\theta_{old}} \gets \pi_{\theta}$
      \State For every question $q$ in $\mathcal{D}_b$, sample $G$ outputs $\{o_i\}_{i=1}^G \sim \pi_{\theta_{old}}(\cdot | q)$
      \State Evaluate rewards $\{r_i\}_{i=1}^G$ for each output $o_i$ using $r_{\varphi}$
      \State For each $o_i$, estimate the group relative advantage $\hat{A}_{i,t}$ at the $t$-th token
      \For{GRPO iteration = 1, \ldots, $\mu$}
        \State Update policy parameters $\pi_\theta$ to maximize the GRPO loss (see Equation \ref{eq:GC-GRPO})
      \EndFor
    \EndFor 
    \State Perform continual learning of $r_\varphi$ with replay strategy.
    \EndFor 
  \end{algorithmic}
  \textbf{Output} $\pi_\theta$
  \label{alg:iter-grpo}
\end{algorithm}

\subsubsection{Outcome Supervision RL with GRPO} Specifically, for a given question $q$, a collection of responses $\{o_1, o_2, \cdots, o_G\}$ is generated by sampling from the previous policy $\pi_{\theta_{old}}$. Each response is then evaluated using a reward model, resulting in a set of rewards $\mathbf{r} = \{r_1, r_2, \cdots, r_G\}$. The obtained rewards undergo normalization, which involves subtracting the mean of the group and dividing by their standard deviation. Within outcome supervision, the normalized reward is allocated to the end of each sequence $o_i$ and assigned as the advantage $\hat{A}_{i, t}$ for all tokens within that sequence, formally written as $\hat{A}_{i, t} = \widetilde{r}_i = \frac{r_i- {\rm mean}(\mathbf{r})}{{\rm std}(\mathbf{r})}$. The policy is then updated to maximize the objective indicated in equation (\ref{eq:GRPO-obj}).

\subsubsection{Process Supervision RL with GRPO}

In outcome supervision, feedback is only provided at the conclusion of each generated output, which may not be sufficiently informative or efficient for guiding policies on intricate mathematical reasoning tasks. Building upon the framework proposed in \cite{wang2023math}, we further investigate process supervision, which assigns rewards after each intermediate reasoning step. More precisely, for a given question $q$ and $G$ sampled solutions $\{o_1, o_2, \cdots, o_G\}$, a process reward model evaluates each step, producing a reward sequence: $\mathbf{R} = \{ \{r_1^{index(1)},\cdots,r_1^{index(K_1)}\}, \cdots,  \{r_G^{index(1)},\cdots,r_G^{index(K_G)}\} \}$. Here, $index(j)$ identifies the end token of the $j$-th step, and $K_i$ corresponds to the number of reasoning steps in the $i$-th generated solution. To ensure comparability, rewards are standardized by subtracting the mean and dividing by the standard deviation of all step-wise rewards: $\widetilde{r}_i^{index(j)} = \frac{r_i^{index(j)} - {\rm mean(\mathbf{R})}}{{\rm std(\mathbf{R})}}$.

With these normalized rewards, process supervision then computes the token-level advantage by aggregating all normalized rewards for subsequent steps, i.e., $\hat{A}_{i, t} = \sum_{index(j) \ge t} \widetilde{r}_i^{index(j)}$. The policy parameters are then updated to maximize the objective function described in equation (\ref{eq:GRPO-obj}).

\subsubsection{Iterative RL with GRPO} During the course of reinforcement learning, the existing reward model may become inadequate for effectively guiding the policy model as it evolves. To address this, we investigate an iterative RL framework using GRPO. As outlined in Algorithm \ref{alg:iter-grpo}, this approach involves constructing updated training datasets for the reward model using samples generated by the current policy model. The reward model is then further refined through continuous training, employing a replay strategy that preserves 10\% of previous data to maintain stability. Subsequently, the policy model is designated as the new reference model, and optimization of the policy continues utilizing the refreshed reward model.

\subsection{Training and Evaluating DeepSeekMath-RL}

Our reinforcement learning experiments utilize DeepSeekMath-Instruct 7B as the foundation. The RL training leverages approximately 144K chain-of-thought-formatted questions pertinent to GSM8K and MATH, all sourced from the SFT dataset. To specifically analyze the effects of RL on evaluation sets absent during RL, we purposefully exclude other SFT-derived questions from the training corpus. The reward model's training data are curated as described in \citep{wang2023math}. The initial reward model is built on the DeepSeekMath-Base 7B architecture and trained with a learning rate of $2 \times 10^{-5}$. For GRPO, the learning rate for the policy model is set to $1 \times 10^{-6}$, and a KL coefficient of $0.04$ is applied. Each question undergoes sampling to generate $64$ possible responses. Training uses a maximum sequence length of 1024 and a batch size of 1024. The policy model undergoes a single update subsequent to each exploration phase. For evaluation, we assess DeepSeekMath-RL 7B using the same benchmarks as DeepSeekMath-Instruct 7B. Here, GSM8K and MATH (with chain-of-thought solutions) serve as in-domain benchmarks, whereas all other tasks are categorized as out-of-domain.

Table \ref{tab:sft_rl_math} presents a comparison of open- and closed-source models employing both chain-of-thought and tool-integrated reasoning across English and Chinese benchmark tasks. Our observations are as follows: 1) With chain-of-thought reasoning, DeepSeekMath-RL 7B achieves 88.2\% accuracy on GSM8K and 51.7\% on MATH. These results not only exceed those of all open-source models between 7B and 70B parameters but also outperform most closed-source counterparts. 2) Notably, DeepSeekMath-RL 7B is exclusively trained on chain-of-thought-style instruction tuning datasets derived from GSM8K and MATH, building upon DeepSeekMath-Instruct 7B. Although its training data is relatively limited, DeepSeekMath-RL 7B consistently outperforms DeepSeekMath-Instruct 7B in every evaluated metric, highlighting the substantial benefits of reinforcement learning.

\section{Discussion}
This section summarizes key insights gained from our pre-training and reinforcement learning experiments.

\subsection{Lessons Learnt in Pre-Training}

We begin by detailing our observations from the pre-training phase. Except where otherwise noted, the experimental protocols described in Section \ref{sec:quality-policy} are followed throughout. Importantly, references to the DeepSeekMath Corpus within this section pertain to a dataset comprising 89B tokens, obtained during the second round of data acquisition.

\subsubsection{Code Training Benefits Mathematical Reasoning}

An often-cited but unproven theory claims that engaging in code training enhances reasoning capabilities. In this work, we aim to address this assertion, focusing on mathematics: code training enhances models' proficiency in mathematical reasoning tasks, regardless of whether external tools are utilized.

To investigate the impact of code training on mathematical reasoning, we conducted experiments under both two-stage and one-stage training regimes:

\textbf{Two-Stage Training}

\begin{itemize}[topsep=0pt]
    \item \textbf{Code Training for 400B Tokens $\rightarrow$ Math Training for 150B Tokens}: We pre-train DeepSeek-LLM 1.3B using 400B tokens sourced from code, then continue training for an additional 150B tokens with a focus on mathematics;
    \item \textbf{General Training for 400B Tokens $\rightarrow$ Math Training for 150B Tokens}: To serve as a baseline, we also conduct training with 400B general domain tokens (drawn from a comprehensive general-purpose dataset constructed by DeepSeek-AI) in the initial stage, instead of code tokens, to assess whether code data provides specific benefits for mathematical reasoning compared to broad-domain textual data.
\end{itemize}

\textbf{One-Stage Training}

\begin{itemize}[topsep=0pt]
    \item \textbf{Math Training Using 150B Tokens}: The DeepSeek-LLM 1.3B model undergoes training solely on 150B math tokens;
    \item \textbf{Joint Training with 400B Code Tokens and 150B Math Tokens}: Sequentially fine-tuning on math data after code data leads to diminished coding abilities. We thus assess whether interleaving code and math tokens in a unified training phase can both foster better mathematical reasoning skills and reduce the extent of catastrophic forgetting.
\end{itemize}

\paragraph{Results}
Table \ref{tab:code-math} and Table \ref{tab:code-math-general-eval} report the results on downstream tasks for each respective training strategy.

Training with code contributes positively to program-aided mathematical reasoning in both the two-stage and one-stage training paradigms. As illustrated in Table \ref{tab:code-math}, code-focused training alone leads to marked improvements in solving GSM8K and MATH tasks via Python in the two-stage scenario. Introducing mathematical reasoning training in the subsequent stage leads to further performance gains. Notably, when adopting the one-stage approach, jointly incorporating code and mathematical tokens alleviates the catastrophic forgetting commonly observed in two-stage training. Moreover, this integrated strategy not only benefits general coding abilities (Table \ref{tab:code-math-general-eval}) but also enhances program-aided mathematical reasoning capabilities (Table \ref{tab:code-math}).

Training on code also enhances mathematical reasoning abilities in contexts where tools are not utilized. In a two-stage training protocol, the preliminary phase focused on code brings about notable, albeit moderate, improvements. Moreover, this code-centric pretraining increases the effectiveness of later mathematical instruction, culminating in superior overall results. In contrast, a one-stage approach that integrates both code tokens and math tokens appears to impair mathematical reasoning performance when no tools are present. This may be attributed to the limited capacity of DeepSeek-LLM 1.3B, which seems insufficient to comprehensively absorb both code and mathematical datasets at the same time.

\subsubsection{ArXiv Papers Seem Ineffective in Improving Mathematical Reasoning}

ArXiv papers are frequently utilized as part of mathematical pre-training datasets \citep{minerva,gpt-f,llemma,mathpile}. Nonetheless, there has been limited in-depth exploration of how their inclusion influences mathematical reasoning abilities. Somewhat surprisingly, our empirical findings suggest that arXiv papers offer little benefit in enhancing mathematical reasoning skills. We perform evaluations across a range of model capacities, notably DeepSeek-LLM 1.3B and DeepSeek-Coder-Base-v1.5 7B \citep{deepseek-coder}, employing arXiv-based datasets subjected to different data processing strategies:

\begin{itemize}[topsep=0pt]
    \item \textbf{MathPile} \citep{mathpile}: This is an 8.9B-token dataset, curated through cleaning procedures and heuristic filtering, with scientific arXiv papers comprising over 85\% of its content.
    \item \textbf{ArXiv-RedPajama} \citep{redpajama}: This dataset consists of all arXiv LaTeX files, with preambles, comments, macros, and bibliographies excluded, yielding a total of 28.0B tokens.
\end{itemize}

In our study, we independently train DeepSeek-LLM 1.3B for 150B tokens and DeepSeek-Coder-Base-v1.5 7B for 40B tokens on each respective arXiv dataset. The results indicate that exposure to arXiv papers alone does not enhance models’ abilities in mathematical reasoning. Specifically, when trained exclusively on the arXiv corpus, neither model shows significant gains—sometimes even declining performance—on a range of mathematical assessment tasks spanning varying levels of difficulty as considered in this research. The evaluations include quantitative reasoning tests such as GSM8K and MATH (Table \ref{tab:arxiv-cot}), multiple-choice exams such as MMLU-STEM (Table \ref{tab:arxiv-cot}), as well as formal mathematical problems as found in miniF2F (Table \ref{tab:arxiv_atp}).

Nonetheless, these findings should be interpreted cautiously due to several limitations. In particular, we have yet to investigate:

\begin{itemize}[topsep=0pt]
    \item How arXiv tokens influence certain mathematical tasks beyond the scope of this study, for instance, the informalization of theorems, which involves translating formal statements or proofs into their informal counterparts;
    \item The outcomes resulting from integrating arXiv tokens with additional categories of data;
    \item If the advantages conferred by arXiv papers persist or change when scaling up to larger models.
\end{itemize}

Therefore, additional investigation is necessary and is beyond the scope of this work.

\subsection{Insights of Reinforcement Learning}
\subsubsection{Towards to a Unified Paradigm} Here, we introduce a cohesive framework that facilitates the analysis of various training algorithms, including SFT, RFT, DPO, PPO, and GRPO. We also present experimental results that probe the key components within this unified framework. Typically, the gradient of a training algorithm with respect to its parameter $\theta$ is expressed as:

\begin{equation}
    \nabla_{\theta}\mathcal{J}_{\textcolor{red}{\mathcal{A}}}(\theta) = \mathbb{E}[\underbrace{(q,o) \sim \textcolor{red}{\mathcal{D}}}_{Data \ Source}]\left( \frac{1}{|o|} \sum_{t=1}^{|o|}  \underbrace{GC_{{\mathcal{A}}}(q, o, t, \textcolor{red}{\pi_{{rf}}})}_{Gradient \ Coefficient}  \nabla_{\theta}\log \pi_{\theta}(o_t | q, o_{<t})\right).
\label{eq:objective}
\end{equation}

This framework comprises three principal elements: (1) the \textit{Data Source $\mathcal{D}$}, which specifies the collection of training instances; (2) the \textit{Reward Function $\pi_{{rf}}$}, responsible for generating the feedback signal used during training; and (3) the \textit{Algorithm $\mathcal{A}$}, which utilizes both the training data and reward signals to produce the gradient coefficient $GC$ that scales the amount of reinforcement or penalty assigned to each sample. We discuss a range of well-known approaches within this generalized structure:

\begin{itemize}
    \item  \textbf{Supervised Fine-tuning (SFT)}: SFT adapts a pretrained model to a specific task by training it further on a dataset curated by humans.
    \item \textbf{Rejection Sampling Fine-tuning (RFT)}: RFT continues training the SFT-tuned model using a subset of its generated responses that pass a filtering process based on correctness in response to SFT dataset prompts.
    \item \textbf{Direct Preference Optimization (DPO)}: DPO enhances the SFT model by training it on additional outputs generated from the SFT model itself, utilizing a pairwise DPO loss objective.
    \item \textbf{Online Rejection Sampling Fine-tuning (Online RFT)}: In contrast to traditional RFT, Online RFT begins with the SFT model as a starting point and incrementally updates the policy model by fine-tuning on augmented samples produced dynamically by the current policy.
    \item \textbf{PPO/GRPO}: PPO/GRPO start with the SFT model to initialize the policy and apply reinforcement learning, using outputs sampled from the up-to-date policy during training.
\end{itemize}

A summary of the components of these approaches is presented in Table \ref{tab:data_policy}. For a comprehensive explanation of the derivation procedure, see Appendix \ref{app:analysis-rl}.

\paragraph{Observation about Data Source} We categorize the data source into two distinct types: online sampling and offline sampling. Online sampling refers to acquiring training data directly from the active exploration of the currently updated policy model, whereas offline sampling utilizes data generated by the initial SFT model's sampling process. Methods such as RFT and DPO are characterized by offline sampling, in contrast to Online RFT and GRPO, which employ online sampling.

As illustrated in Figure \ref{fig:rl-analysis}, our results indicate that Online RFT substantially surpasses RFT across both evaluated benchmarks. In the initial phases of training, Online RFT and RFT display similar performance, but as training progresses, Online RFT establishes a clear lead, highlighting the benefits of online training strategies. This outcome is expected: during early training, the actor and the SFT model are quite similar, resulting in sampled data that reflect only slight deviations. In contrast, during later training stages, the data drawn from the actor diverge more prominently, thereby making real-time data sampling markedly more beneficial.

\paragraph{Observation about Gradient Coefficient} The algorithm leverages the reward signal to determine the gradient coefficient, which in turn guides the updates of the model parameters. In our experimental setup, we categorize the reward function into two components: ‘Rule’ and ‘Model’. The ‘Rule’ component assesses the response by verifying the correctness of the answer, whereas the ‘Model’ involves training a dedicated reward model that assigns a score to each response. Notably, the reward model’s training data is sourced from the judgments made by the ‘Rule’ criteria. As shown in Equations \ref{eq:GC-RFT} and \ref{eq:GC-GRPO}, GRPO and Online RFT differ fundamentally: GRPO modulates its gradient coefficient dynamically in relation to the reward score derived from the reward model. This mechanism enables nuanced reinforcement or penalization, scaling with the magnitude of each response’s reward. Conversely, Online RFT does not possess this capability; it applies uniform reinforcement to all correctly answered responses and refrains from penalizing incorrect responses altogether.

As illustrated in Figure \ref{fig:rl-analysis}, GRPO outperforms online RFT, emphasizing the effectiveness of modifying both positive and negative gradient coefficients. Moreover, GRPO+PS achieves better results than GRPO+OS, which demonstrates the advantage of employing more detailed, step-specific gradient coefficients. Additionally, we investigate the use of iterative RL; in our setup, two iterations are carried out. Figure \ref{fig:iter-rl} reveals that iterative RL markedly enhances performance, with the most substantial gains observed during the initial iteration.

\subsubsection{Why RL Works?} In this study, we apply reinforcement learning to a specific subset of instruction tuning data and observe marked gains over models solely utilizing instruction tuning. To elucidate the underlying reasons for the effectiveness of reinforcement learning, we assess both Pass@K and Maj@K accuracies for the Instruct and RL models using two different benchmarks. Figure \ref{fig:maj-pass} demonstrates that, while RL does not increase Pass@K accuracy, it does yield improvements in Maj@K. This suggests that reinforcement learning enhances general model reliability by fortifying the robustness of the output distribution—specifically, \textbf{the improvement appears to stem from increasing the likelihood of producing a correct answer among the TopK responses, rather than elevating the model’s base-level skill}. Analogously, \citep{wang2023making} point out a \textbf{misalignment problem} in SFT models during reasoning tasks, highlighting that a variety of preference alignment techniques can lead to better reasoning results for SFT models \citep{yuan2023rrhf,song2023preference,wang2023making}.

\subsubsection{How to Achieve More Effective RL?}
\label{sec:effec_rl}
Our results show that RL can be highly effective in tasks involving mathematical reasoning. Additionally, we introduce a comprehensive framework to interpret several prominent training strategies. In this unified view, these approaches are regarded as variants of either full or approximated RL methodologies. Equation \ref{eq:objective} outlines three crucial elements: Data Source, Algorithm, and Reward Function. We also suggest avenues for future exploration concerning these components.

\paragraph{Data Source} The data source serves as the foundational input for all training methodologies. In RL settings, we define the data source as comprising unlabeled questions, along with responses generated from the policy model. For this study, our dataset is restricted to questions utilized during the instruction tuning phase, and output responses are obtained using a basic nucleus sampling strategy. This limitation may contribute to our RL pipeline enhancing only the Maj@K metric. Looking ahead, we plan to extend our RL framework to accommodate out-of-distribution prompts and to employ \textbf{advanced sampling (decoding) strategies}, such as those employing tree-search algorithms \citep{yao2023tree}. Moreover, incorporating \textbf{efficient inference techniques} \citep{xia-etal-2023-speculative,leviathan2023fast,kwon2023efficient,xia2024unlocking}, which crucially impact how effectively policy models explore, will be a key focus in future work.

\paragraph{Algorithms} Algorithms utilize both the data and the reward signal to compute the gradient coefficient for updating model parameters. According to Equation \ref{eq:objective}, contemporary approaches generally place complete \textbf{TRUST} in the reward function, using its signal to modulate the conditional probability of specific tokens. Nevertheless, the reliability of the reward signal cannot be universally guaranteed, particularly when tackling highly intricate tasks. As an illustration, the PRM800K datasets \citep{lightman2023let}, despite rigorous annotation efforts by skilled annotators, still exhibit an error rate of roughly 20\% in their annotations\footnote{\url{https://github.com/openai/prm800k/issues/12\#issuecomment-1728491852}}. Consequently, our focus will shift to reinforcement learning algorithms that demonstrate robustness to noisy reward signals. We anticipate that such \textbf{WEAK-TO-STRONG} \citep{burns2023weak} alignment methodologies have the potential to fundamentally transform existing learning algorithms.

\paragraph{Reward Function} The reward function serves as the primary training signal. In RL settings, this function is typically instantiated by a neural reward model. There are three significant research directions concerning reward models: 1) \textbf{Improving the generalization capability of the reward model.} The reward model needs to generalize effectively to tackle out-of-distribution queries and advanced generation outputs. If it fails in this aspect, reinforcement learning might only reinforce the existing output distribution of LLMs, rather than boosting their underlying performance; 2) \textbf{Capturing and expressing uncertainty in the reward model.} Accurately representing uncertainty may serve as a critical connector between limited reward models and weak-to-strong learning frameworks; 3) \textbf{Developing high-quality, efficient process reward models} that deliver detailed, step-level feedback during reasoning \citep{lightman2023let,wang2023math}.

\section{Conclusion, Limitation, and Future Work}

In this work, we introduce DeepSeekMath, which surpasses all other open-source models on the MATH benchmark at the competition level, and its performance closely rivals that of proprietary models. DeepSeekMath is based on the DeepSeek-Coder-v1.5 7B architecture and is further trained for 500B tokens, of which 120B are mathematics-focused tokens derived primarily from Common Crawl. Our comprehensive ablation analysis demonstrates that web pages are a particularly valuable source of high-quality mathematical training material, whereas data from arXiv does not contribute as much as anticipated. We propose Group Relative Policy Optimization (GRPO), an adaptation of Proximal Policy Optimization (PPO), which enhances mathematical reasoning performance while reducing memory requirements. Experimental findings confirm that GRPO remains effective even when DeepSeekMath-Instruct 7B already performs strongly on standardized evaluations. Additionally, we offer a unified framework for interpreting related methodologies and outline several promising directions for advancing reinforcement learning efficacy.

While DeepSeekMath~demonstrates strong performance on benchmarks involving quantitative reasoning, its abilities in geometry and theorem proving are somewhat inferior to those of proprietary models. For example, our preliminary tests revealed that the model struggles with triangle and ellipse-related questions, which may suggest that both pre-training and fine-tuning stages suffered from selection bias in data curation. Moreover, due to its limited model size, DeepSeekMath~lags behind GPT-4 in few-shot learning scenarios. GPT-4 is able to boost its accuracy with a small number of exemplars, whereas DeepSeekMath~shows minimal difference between its zero-shot and few-shot results. Looking ahead, we plan to enhance our engineered data selection process to generate a more robust and diverse pre-training corpus. We also aim to investigate potential avenues for more efficient reinforcement learning of LLMs, as outlined in Section \ref{sec:effec_rl}.

\end{document}